% Part: sets-functions-relations
% Chapter: relations
% Section: relations-as-sets

\documentclass[../../../include/open-logic-section]{subfiles}

\begin{document}

\olfileid{sfr}{rel}{set}
\olsection{Relaties als verzamelingen}

\begin{explain}
In \olref[sfr][set][imp]{sec} noemden we een paar belangrijke
verzamelingen: $\Nat$, $\Int$, $\Rat$, $\Real$. Je kent vast nog wel
een paar interessante relaties tussen de !!{element}s van sommige van
die verzamelingen. Zo heeft elk van deze verzamelingen een gebruikelijke
manier om de elementen op volgorde te zetten: een \emph{ordening}. Er
is ook de relatie \emph{is hetzelfde als}. Ieder
object heeft die relatie met zichzelf, en met niets anders. We komen
nog veel meer interessante relaties tegen, en er zijn er nog veel meer
denkbaar. Voordat we ze langslopen, laten we eerst zien dat je een
relatie kunt bekijken als een bijzonder soort verzameling.
  
Daarvoor hebben we twee dingen uit \olref[sfr][set][pai]{sec} nodig.
Het eerste is een \emph{geordend paar}: met $a$ en $b$ kun je
het paar~$\tuple{a, b}$ maken. Daarbij doet de volgorde van de elementen
er \emph{wel} toe. Dus als $a \neq b$, dan is $\tuple{a, b} \neq
\tuple{b, a}$. Bij een ongeordend paar is dat anders. Zo'n paar is een
verzameling met $2$ elementen, en daarvoor geldt $\{a, b\}=\{b, a\}$.
Het tweede begrip is het \emph{cartesisch product}. Als $A$ en $B$
verzamelingen zijn, kun je~$A \times B$ maken. Dat is de verzameling
van alle paren $\tuple{x, y}$ waarvoor $x \in A$ en $y \in B$. Dus
$A^{2}= A \times A$ bestaat uit alle geordende paren waarvan beide
elementen uit~$A$ komen.
  
Nu kijken we naar één concrete relatie op een verzameling: de
$<$-relatie op de verzameling~$\Nat$ van natuurlijke getallen. Neem
alle getallenparen $\tuple{n, m}$ waarvoor $n<m$. Dat geeft de
verzameling
\[
  R=\Setabs{\tuple{n, m}}{n, m \in \Nat \text{ en } n<m}.
\]
De bewering dat $n$ kleiner is dan $m$ en de bewering dat
$\tuple{n, m}$ in $R$ zit, zeggen precies hetzelfde:
\[
      n<m\text{ precies wanneer }\tuple{n, m} \in R.
\]
We verliezen dus geen informatie als we zeggen dat de verzameling $R$
zelf de $<$-relatie op $\Nat$ \emph{is}.

Zo kun je bij elke relatie tussen getallen een deelverzameling van
$\Nat^{2}$ maken. Het werkt ook andersom. Elke verzameling
getallenparen $S \subseteq \Nat^{2}$ bepaalt welke relatie erbij
hoort: $n$ heeft die relatie met $m$ precies als $\tuple{n, m} \in S$.
Daarom gebruiken we de volgende definitie:
\end{explain}
  
\begin{defn}[Binaire relatie]
Een \emph{binaire relatie} op een verzameling $A$ is een
deelverzameling van~$A^{2}$. Als $R \subseteq A^{2}$ zo'n binaire
relatie op~$A$ is en $x, y \in A$, gebruiken we soms $Rxy$ (of $xRy$)
als korte schrijfwijze voor $\tuple{x, y} \in R$.
\end{defn}
  
\begin{ex}
  \ollabel{relations}
Je kunt alle paren natuurlijke getallen in $\Nat^{2}$ zo in een tabel
met rijen en kolommen zetten:
\[
  \begin{array}{ccccc}
  \mathbf{\tuple{ 0,0 }} & \tuple{ 0,1 } &
    \tuple{ 0,2 } & \tuple{ 0,3 } & \ldots\\
  \tuple{ 1,0 } & \mathbf{\tuple{ 1,1 }} &
    \tuple{ 1,2 } & \tuple{ 1,3 } & \ldots\\
  \tuple{ 2,0 } & \tuple{ 2,1 } &
    \mathbf{\tuple{ 2,2 }} & \tuple{ 2,3 } & \ldots\\
  \tuple{ 3,0 } & \tuple{ 3,1 } & \tuple{ 3,2 } &
    \mathbf{\tuple{ 3,3 }} & \ldots\\
  \vdots & \vdots & \vdots & \vdots & \mathbf{\ddots}
  \end{array}
\]
We hebben de diagonaal vet gedrukt. De paren op die diagonaal vormen
namelijk de deelverzameling van $\Nat^2$
\[
  \{\tuple{0,0 }, \tuple{ 1,1 }, \tuple{ 2,2 }, \dots\},
  \]
en deze deelverzameling is de \emph{identiteit op}~$\Nat$. (Deze relatie
gebruiken we vaak. Daarom leggen we de notatie
$\Id{A}=\Setabs{\tuple{ x,x }}{x \in A}$ vast voor elke verzameling
$A$.) De paren boven de
diagonaal vormen de deelverzameling
\[
  L = \{\tuple{ 0,1 },\tuple{ 0,2 },\ldots,\tuple{ 1,2 },
  \tuple{ 1,3 }, \dots, \tuple{ 2,3 }, \tuple{ 2,4 },\ldots\},
\]
en dit is de relatie \emph{kleiner dan}: $Lnm$ geldt precies wanneer
$n<m$. De paren onder de diagonaal vormen de deelverzameling
\[
  G=\{ \tuple{ 1,0 },\tuple{ 2,0 },\tuple{
    2,1 }, \tuple{ 3,0 },\tuple{ 3,1 },\tuple{ 3,2 }, \dots\},
\]
en dit is de relatie \emph{groter dan}: $Gnm$ geldt precies wanneer
$n>m$. Neem nu de vereniging van $L$ en de zojuist gedefinieerde
identiteit~$I$, en noem die $K=L\cup I$.
Dit is de relatie \emph{kleiner dan of gelijk aan}: $Knm$ geldt precies
wanneer $n \le m$. Zo is $H=G \cup I$ de relatie \emph{groter dan of
gelijk aan.} De relaties $L$, $G$, $K$ en $H$ zijn van een bijzonder
soort. Ze heten \emph{ordeningen}. $L$ en $G$ hebben allebei deze
eigenschap: $L$ koppelt geen enkel getal aan zichzelf en $G$ ook niet
(dus voor elk $n$ geldt niet $Lnn$ en
ook niet $Gnn$). Een relatie met deze eigenschap heet
\emph{irreflexief}. Als zo'n relatie ook een ordening is, heet ze een
\emph{strikte ordening.}
\end{ex}

\begin{explain}
Ordeningen en identiteit zijn belangrijke relaties die vanzelfsprekend
aanvoelen. Toch is volgens onze definitie \emph{elke} deelverzameling
van $A^{2}$ een relatie op~$A$, hoe vreemd of bedacht die ook lijkt. Zo
is $\emptyset$ een relatie op elke verzameling: de \emph{lege relatie},
die voor geen enkel paar elementen geldt. Ook $A^{2}$~zelf is een
relatie op~$A$. Die geldt voor elk paar en heet de \emph{universele
relatie}. Zelfs iets als
$E=\Setabs{\tuple{n, m}}{n>5 \text{ of } m \times n \ge 34}$ telt als
een relatie.
\end{explain}
  
\begin{prob}
  Schrijf alle !!{element}s op van de relatie $\subseteq$ op de verzameling
  $\Pow{\{a, b, c\}}$.
\end{prob}
\end{document}
